\documentclass[10pt,twocolumn]{article}
\usepackage[margin=0.72in]{geometry}
\usepackage{fontspec}
\setmainfont{Latin Modern Roman}
\usepackage{amsmath,amssymb}
\usepackage{graphicx}
\usepackage{booktabs}
\usepackage{tabularx}
\usepackage{array}
\usepackage{microtype}
\usepackage{enumitem}
\usepackage{caption}
\usepackage[hidelinks]{hyperref}
\usepackage{xurl}
\graphicspath{{../results/figures/campaign_c_20261007/}}
\setcounter{secnumdepth}{3}
\setlength{\columnsep}{0.25in}
\setlength{\tabcolsep}{3pt}
\setlength{\parindent}{1em}
\setlength{\parskip}{0pt}
\setlength{\emergencystretch}{2em}
\captionsetup{font=small,labelfont=bf}
\setlist{nosep,leftmargin=*}
\providecommand{\tightlist}{\setlength{\itemsep}{0pt}\setlength{\parskip}{0pt}}
\title{Transparent Risk-Sensitive RIS Phase Control Under CSI Uncertainty}
\author{%
\begin{tabular}{c}
Pengliang Tian\textsuperscript{1}, Pengyuan Zhang\textsuperscript{1}, Haoshan Sun\textsuperscript{1}, Changjiang Zhang\textsuperscript{1,*}\\[0.35ex]
\parbox{0.95\textwidth}{\centering\footnotesize \textsuperscript{1}Department of Communication and Information Engineering, Century College, Beijing University of Posts and Telecommunications, Beijing, China}\\[0.35ex]
\footnotesize First-author ORCID: 0009-0003-5941-9797; corresponding author: \href{mailto:zhangchangjiang@ccbupt.cn}{zhangchangjiang@ccbupt.cn}, ORCID: 0009-0002-9936-1377
\end{tabular}
}
\date{}
\begin{document}
\maketitle
\begin{abstract}

Reconfigurable intelligent surfaces (RISs) make the propagation
environment a reconfigurable control variable, but a phase action is
only as reliable as the channel information used to select it. We study
a causal controller for a multi-user normalized RIS simulator in which
channel-estimation error drifts over time. The controller constructs a
finite, posterior-like ensemble of error scenarios, selects a quantized
phase action using a per-user lower-tail feasibility check, and reuses
the action until a fixed trigger requests a refresh. Candidate search,
validation, phase quantization, refusal handling, and phase-control
rate-evaluation cost are explicit.

The frozen held-out Campaign C contains 20,480 rows over 16 factor
conditions, ten independent seed blocks, four episodes, eight slots, and
four controllers. Using paired bootstrap over complete seed blocks, the
maximum one-sided 95\% upper bound for absolute weak-user outage is
0.000; against the always-risk reference, the outage-delta upper bound
is 0.000, the goodput relative lower 05 bound is −0.00615, and the
phase-control rate-evaluation reduction lower 05 bound is 0.33398. These
results support a conditional reliability--computation trade-off on the
declared normalized simulator grid. They do not establish a calibrated
population posterior, a learned or reinforcement-learning policy,
physical phase-write timing, OTA performance, or universal superiority.
The Jetson timing record is simulator diagnostic data; no RIS
phase-write endpoint or physical slot trace was available.

\end{abstract}

\textbf{Keywords:} reconfigurable intelligent surface; CSI uncertainty;
temporal correlation; phase reuse; weak-user outage; risk-aware control;
paired seed-block evaluation.

\paragraph{Highlights}

\begin{itemize}
\tightlist
\item
  A causal per-user guard makes RIS phase reuse and lower-tail
  feasibility explicit.
\item
  A frozen 20,480-row Campaign C passes the four registered always-risk
  simulator gates.
\item
  Physical phase-write timing is kept as UNKNOWN rather than inferred
  from CUDA timing.
\end{itemize}

\section{Introduction}\label{introduction}

RIS-assisted communication uses a programmable reflecting surface to
shape the effective channel between a transmitter and one or more users
{[}2,6,11,13{]}. The phase profile is attractive as a control variable
because it can be changed without adding a full radio-frequency chain at
every element. At the same time, the phase profile is selected from
information that is neither exact nor static. Pilot noise, model
mismatch, channel tracking, and channel aging all affect the estimate
available at a later payload slot {[}14,22--24,27,30--32{]}. In a
multi-user link, a phase vector that improves a mean or sum rate can
still leave the weakest user close to an outage threshold, a concern
addressed by robust RIS beamforming and outage-constrained designs
{[}7,8,21,28,29{]}. The resulting engineering problem is therefore a
sequential control problem under uncertain and correlated information
rather than a one-shot point-estimate optimization.

The most direct response to uncertainty is to refine the channel
estimate and recompute an action frequently. Tracking, beam prediction,
and two-timescale estimation reduce this burden in different ways
{[}3,4,14,27,30{]}, but refinement still has a measurable cost. It can
increase pilot and estimation work, candidate generation, rate
evaluations, host-device transfers, and the time between the first
available observation and the phase write used by the data slot. Reusing
a nominal action has the opposite profile: it can be inexpensive when
the channel is stable, but it has no explicit mechanism for deciding
whether the uncertainty around the reused action is acceptable.
Learning-based and reinforcement-learning controllers provide another
design route {[}5,9,12, 19,25,31,33{]}, but their training and policy
state make a matched causal cost audit a different estimand. A useful
controller must expose this reliability and work trade-off instead of
hiding it in an opaque policy or in an average rate alone.

We study a transparent mechanism for this purpose. Given a current
channel estimate, the controller forms a finite ensemble of plausible
CSI errors. It evaluates candidate quantized RIS phase actions over that
ensemble and keeps the candidates whose lower-tail service statistic
passes a configured feasibility rule. A trigger compares the current
reused action with the available uncertainty evidence and requests a
refresh when reuse is no longer acceptable. The scenario construction,
lower-tail check, trigger, phase quantization, refusal path, and cost
accounting are explicit. This makes the mechanism inspectable and
permits a matched comparison with periodic risk reuse, always-refine
control, and periodic nominal-CSI reuse.

The evidence is separated into a completed statistical campaign and a
hardware boundary. The corrected held-out Campaign C is complete and
hash-linked: it contains 20,480 rows, 16 factor conditions, 64 condition
summaries, and 640 paired seed/controller streams. The four primary
statistical thresholds were frozen before the held-out seeds were
inspected. All four pass against the always-risk reference. Comparisons
with periodic-risk and periodic-nominal baselines retain their service
results but fail the all-condition work-reduction gate because those
baselines use different refresh schedules; this is reported as a
sensitivity result rather than silently reclassified as a success. The
first erroneous launch remains a 2,450/5,120-row interrupted artifact
and is excluded from inference.

The physical timing claim is intentionally unresolved. The Jetson
exposes generic SPI/GPIO device nodes but no detected RIS controller or
phase-commit endpoint. Consequently, the simulated total-latency P95
(maximum 388.083 ms) is a diagnostic for algorithmic cost only. It is
not a phase-write completion measurement and cannot be compared with a
physical slot. The manuscript uses the results as a normalized RIS-only
simulation study and keeps OTA, real-time, deployment, and
calibrated-population language outside the claim set.

This study makes four bounded contributions. First, it gives an
inspectable causal phase-reuse mechanism in which the CSI-error
ensemble, lower-tail rule, guard, trigger, quantization, refusal path,
and cost ledger are explicit. Second, it contributes a frozen
development/held-out protocol with complete seed-block paired bootstrap
rather than slot-level pseudo-replication. Third, it reports absolute
all-slot weak-user service together with relative goodput and online
work, making refused or failed actions visible in the denominator.
Fourth, it preserves a clean evidence boundary: the statistical gates
are closed for the declared simulator grid, while physical phase-write
timing and deployment evidence remain unknown. No individual ingredient
is presented as new without the final full-text literature comparison.

The controller and provenance path are summarized in Figure 1.

\section{System Model}\label{system-model}

\subsection{Multi-user RIS link}\label{multi-user-ris-link}

Consider a transmitter with a fixed set of antennas serving a set of
users through a programmable RIS path. Let the RIS contain \(N\)
elements and let its phase vector be
\(\boldsymbol{\phi}=[\phi_1,\ldots,\phi_N]^\mathsf{T}\). Each phase
belongs to the configured discrete alphabet. The effective channel for
user \(u\) is the normalized cascaded transmitter-RIS-user contribution.
The simulator uses normalized channel coefficients and a versioned
configuration for path gains, noise, transmit power, pilot observations,
and the phase alphabet. These quantities define a controlled evaluation
environment; they are not measurements of a particular deployed surface.

For a candidate phase vector and a channel realization, the simulator
computes a user rate from the configured signal and interference model.
We write this quantity as \(R_u(\boldsymbol{\phi},H)\), where \(H\)
denotes the channel state available for evaluation. The controller can
evaluate a candidate on the nominal estimate and on a finite set of
error-perturbed scenarios. A payload slot is successful for user \(u\)
only when its transmitted payload is delivered and its measured or
simulated rate meets the configured outage threshold. The weak-user
index is a condition-level reference estimand: before forming any
controller deltas, it is frozen as the argmin of the always-risk
true-channel mean rate over the matched condition trace. The same index
is then applied to all controllers. This rule is controller-independent
but does not assert a fixed physical user identity across deployments.

The action space includes the RIS phase vector and the associated commit
or refusal decision. Candidate generation and rate evaluation are
explicit operations. Phase quantization is applied before an action is
committed, and the same action space is used by the triggered controller
and every baseline. This prevents an apparent reliability or work
advantage from arising solely from a different phase alphabet or
stopping rule.

\subsection{CSI uncertainty and temporal
evolution}\label{csi-uncertainty-and-temporal-evolution}

At decision time \(t\), the controller receives an observation-derived
channel estimate \(\widehat H_t\). The unknown channel is represented as
an estimate plus a structured error whose scale and temporal correlation
are configuration factors. The error process is generated consistently
for matched controller runs, so that a change in a controller's result
is not caused by a different channel trace. The full campaign uses two
recorded error scales and two recorded temporal-correlation values.
These factors are simulation parameters within the declared grid; they
do not establish behavior outside that grid.

The controller's scenario ensemble is conditional on the current
estimate and configured uncertainty model. It is called posterior-like
because it provides a distribution-shaped collection of plausible errors
for action validation. The term does not assert that the ensemble is
calibrated to a population of physical channels. Calibration, if needed,
would require independent channel measurements and a separate validation
design. Posterior scenario draws are also not physical-channel repeats.
They are nested within an action decision and remain inside the
corresponding independent seed block for uncertainty analysis.

Temporal correlation matters because action reuse spans multiple slots.
A highly correlated process can make reuse appear efficient over a short
window, whereas a less correlated process can invalidate the same action
sooner. The experiment therefore records the correlation factor and
reports condition-level patterns instead of treating slots as
independent samples. Any inference remains conditional on the Campaign C
held-out error/rho grid and does not extend to arbitrary shifts.

\subsection{Service, goodput, and cost
definitions}\label{service-goodput-and-cost-definitions}

For a scheduled slot, define \(served_{u,t}=1\) only when a payload for
user \(u\) is transmitted and the resulting rate meets the configured
threshold. If the controller refuses an action, times out, produces no
decision, fails to write the phase, or has a missing record, then
\(served_{u,t}=0\) and the delivered payload is zero. These events
remain in the denominator. For condition (c), the weak-user all-slot
outage is

\[
O_w(c)=1-\frac{\sum_{e,t} served_{u^*,e,t}}{N_{slots}(c)},
\]

where \(N_{slots}(c)\) includes every scheduled slot. The target is
\(O_w(c)\leq 0.05\) with a one-sided 95\% upper seed-block bound no
greater than 0.05 for every condition. Campaign C closes this simulator
gate with a maximum one-sided upper95 of 0.000; this is not a deployed
service guarantee.

Goodput is the sum of delivered payload bits. A refusal contributes zero
bits, even if a previous action would have produced a good rate. This
accounting prevents a controller from looking reliable by reporting only
the slots it admitted. Relative comparisons, such as a goodput ratio or
an outage difference against a reference, cannot waive the absolute
target.

Online work is recorded in a common ledger. Estimation and posterior
update, candidate generation and search, action validation,
phase-control rate evaluations, transfers and synchronization, and phase
writing each have an explicit boundary. Initialization, including the
first estimate, first phase construction, allocation, or compilation, is
recorded separately and included in campaign totals. The same boundaries
apply to the triggered controller, periodic risk reuse, always-refine
control, and nominal-CSI reuse.

\subsection{Physical timing
boundary}\label{physical-timing-boundary}

Let \(T_{slot}\) be the physical duration of the exact waveform and
pilot schedule used by the communication testbed. Decision latency is
measured from the first input sample available to the controller through
completion of the phase/action write used by that slot. The interval is
decomposed into

\[
L_{dec}=L_{est}+L_{search}+L_{validate}+L_{commit},
\]

where the final term includes host-device transfers, synchronization,
and the phase write. The gate is \(P95[L_{dec}]\leq0.10T_{slot}\), with
a one-sided 95\% upper bound also below the threshold and no scheduled
slot exceeding \(T_{slot}\). A monotonic timer and accelerator-event
synchronization are required. The simulator's \texttt{decision\_ms}
field does not contain this physical slot definition and is diagnostic
only. Every real-time or hardware-feasibility statement remains UNKNOWN
until a physical phase-write and slot trace are measured.

\section{Method}\label{method}

\subsection{Scenario ensemble}\label{scenario-ensemble}

The controller begins with the current estimate and the configured
CSI-error model. It generates a finite ensemble of error scenarios and
combines each scenario with the estimate to obtain a plausible channel
state. The ensemble is generated from controller inputs and configured
uncertainty; future truth is not used to construct it. The number of
scenarios, randomization seed, and error scale are recorded in the
effective configuration. Scenario generation is part of
estimation/posterior work in the cost ledger.

For each candidate phase vector, the controller computes the user-rate
statistics over the scenario ensemble. A lower-tail statistic is used to
represent weak-user risk. A candidate is feasible only if its configured
lower-tail service criterion passes. Among feasible candidates, the
controller applies a deterministic objective and tie-break order. If no
candidate passes, the action is refused or the configured fallback is
applied and logged. The fallback cannot silently remove a service
failure from the all-slot denominator.

The finite ensemble is a decision device rather than a claim of an exact
chance constraint. Its reliability interpretation is conditional on the
scenario model, candidate budget, and simulator. A future physical study
would need to calibrate the uncertainty model and validate the
lower-tail criterion with measured channels. That work is outside the
current evidence.

\subsection{Trigger and action
reuse}\label{trigger-and-action-reuse}

After a phase action is committed, the controller may reuse it for later
slots. The trigger examines the new observation and the current ensemble
evidence. If the action remains feasible under the lower-tail rule,
reuse avoids a full refresh. If the evidence crosses the configured
trigger boundary, the controller generates and validates a new candidate
set. The trigger state, initial action, refresh decision, and commit or
refusal reason are recorded for each slot.

Refresh work is not assumed to save pilots or transfers. The ledger
records estimation, search, validation, rate evaluations, and commit
work when they occur, including the first action. A lower refresh count
can therefore be interpreted only together with measured rate-evaluation
and service outcomes. The trigger does not learn a policy from held-out
traces; it applies the versioned deterministic rule in the effective
configuration.

\subsection{Baselines and matched action
accounting}\label{baselines-and-matched-action-accounting}

The comparison includes posterior-triggered control, periodic risk
reuse, always-refine control, and periodic nominal-CSI reuse. All
controllers use the same candidate/action space, phase quantization,
pilot budget, outage threshold, and stopping rules. Always-refine
recomputes the configured risk-aware action at each scheduled refresh.
Periodic risk reuse refreshes on a fixed period with a frozen risk
margin. Periodic nominal reuse starts from a fixed nominal-CSI action
and refreshes on a fixed period. Initialization is identical in
accounting, even when the controller later reuses an action.

Baseline periods, risk margins, initial actions, and tie-break rules
must be tuned on a development seed set disjoint from held-out
evaluation. The selection rule first requires the upper weak-user outage
bound to satisfy the absolute target, then chooses the highest all-slot
goodput, breaking ties by lower rate-evaluation count and lower P95
latency. A frozen manifest records the selected settings before any
held-out shift seed is read. Existing settings must not be described as
optimal until this protocol is complete.

\subsection{Implementation and
reproducibility}\label{implementation-and-reproducibility}

The implementation exposes source, configuration, effective
configuration, random seed, raw row, summary, and analysis hashes. The
phase quantization and action accounting are verified by unit tests and
raw recomputation. A backend parity check compares the recorded NumPy
and CuPy paths under the same host inputs; parity is an implementation
check, not a deployment or acceleration claim. Every Python entry point
accepts an explicit seed. Configurations live under \path{configs/},
and result manifests identify the configuration hash.

The method is therefore inspectable at three levels: the algorithm code
and unit tests, the per-slot records and action reasons, and the paired
analysis that reconstructs condition-level comparisons. A missing
artifact is not replaced by a plotted estimate. The figure plan and
evidence matrix define the status of each claim before any caption is
finalized.

\subsection{Causal state update and guard
decision}\label{causal-state-update-and-guard-decision}

The temporal state is an augmented complex Gaussian belief over a latent
cascaded channel and its estimation error. At the first slot, the
observed coefficient power initializes the channel and error covariance
using the configured relative error scale. At later slots, the error
covariance is propagated with the configured AR(1) coefficient \(\rho\),
while the channel state follows its declared stationary transition. The
update consumes only the current observation and declared covariance
metadata. The simulator truth and future error realization are never
passed to the controller. The normalized innovation statistic and
relative posterior uncertainty are emitted as audit fields and are used
by the trigger.

On a refresh, the posterior mean supplies the decision-side channel
estimate. The scenario scale is the clipped product of posterior
uncertainty and the frozen inflation factor 1.5. Four scenarios are used
by the full candidate search in Campaign C. The optimizer evaluates the
same 2-bit phase alphabet for every controller, keeps a candidate only
if the configured lower-tail per-user target passes, and applies
deterministic tie-breaking to the expected sum-rate objective. The
chosen phase vector is quantized before it is stored in the action
record. The selected action therefore has an explicit provenance from
observation, posterior state, candidate budget, and quantized phase.

Between refreshes, the retained phase vector is checked by a causal
one-scenario guard. In the frozen guard1 revision, the guard budget and
guard scenario count are both one, with the same per-user target and the
same inflated uncertainty scale as the full search. A passing guard
consumes only its recorded guard evaluation and keeps the action. A
failed guard immediately invokes the full four-scenario refresh; the row
records both the failed guard and the resulting full-refresh work. This
separation makes the work reduction auditable and prevents a cheap guard
from silently changing the phase vector without a logged refresh.

The trigger also enforces a minimum refresh gap and a periodic refresh
horizon. The selected held-out policy has a refresh period of four
slots, an uncertainty threshold of 0.35, an innovation threshold of 4.0,
and no adaptive guard escalation. These values were frozen from the
development split. Adaptive guard code is tested and retained as a
development candidate, but it is not the Campaign C policy; this avoids
choosing a more favorable trigger after reading held-out rows.

\subsection{Action, service, and work
records}\label{action-service-and-work-records}

Each raw row identifies RIS size, user count, error scale, temporal
correlation, seed, episode, slot, controller, and a stream hash. It
stores the selected phase vector, per-user rates, weakest-user index,
outage/service flags, refresh and trigger reasons, posterior
uncertainty, normalized innovation, reuse age, candidate count, rate
evaluations, guard evaluations, unused budget, and timing components.
The row also carries source, configuration, effective configuration,
base-seed, and backend metadata. These fields allow a verifier to
recompute the summary without relying on a plot or a manually
transcribed number.

The cost ledger distinguishes phase-control rate evaluations from
simulator truth evaluation. The latter is used to score an action after
the decision and is excluded from the controller's online work metric.
Initialization is kept in the stream rather than amortized away. For
each baseline, the same action space and phase quantizer are used, and
the same service denominator is applied. This matters because a periodic
baseline can look cheap or expensive solely because it refreshes on a
different schedule. The primary work comparison is therefore registered
against always-risk, whose full risk action is recomputed at each
scheduled refresh; periodic results are shown as sensitivity checks.

The analysis code first validates row uniqueness and condition
completeness, then aggregates within seed blocks, forms paired
controller deltas, and only then applies the bootstrap. The result is a
reproducible chain from raw slot record to condition-level interval. No
posterior scenario draw is promoted to an independent sample, and no
slot is counted as a separate seed replicate.

\subsection{Estimands and uncertainty
reporting}\label{estimands-and-uncertainty-reporting}

The primary absolute estimand is the weak-user all-slot outage for one
condition and controller, defined as the number of weak-user service
failures over every scheduled slot. The reported point estimate is
obtained after within-seed aggregation, and the uncertainty interval
resamples complete seed blocks. This preserves the temporal and episode
structure of a block. A zero observed failure count is not interpreted
as a proof of zero population risk; it simply produces the recorded
one-sided bootstrap bound for the registered finite campaign. The
weak-user index is not selected separately for each controller: for each
factor condition, it is the argmin of the always-risk true-channel mean
rate over the matched condition trace, frozen before paired deltas are
formed. This defines a reproducible condition-level reference estimand,
rather than a fixed physical user identity.

For a target controller and a baseline on the same stream, the outage
delta is the target outage minus the baseline outage. Goodput relative
delta is computed with the baseline goodput as the denominator, with
zero-denominator cases recorded explicitly. The rate-evaluation
reduction is one minus the target phase-control evaluation count divided
by the matched baseline count. All three quantities are formed at the
seed-block level before the paired bootstrap, so the comparison retains
common channel and observation realizations. The primary reference is
always-risk because its full risk action is recomputed at each scheduled
refresh and supplies a stable high-work denominator.

The one-sided limits are the 95th or 5th percentile of the recorded
bootstrap distribution, as appropriate. The analysis seed, bootstrap
draw count, input hashes, and effective configuration hash are written
beside every summary and paired row. The 10,000 draws improve numerical
resolution of the reported finite-sample resampling distribution, but
they do not create additional independent channel blocks. Independent
evidence comes only from the ten predeclared held-out seeds.

As a non-selection sensitivity audit, the release also reports a
one-sided Clopper--Pearson bound for the weak-user failure count over
the scheduled slots. This exact bound is not substituted for the
registered seed-block bootstrap gate and does not create additional
independent samples; it provides a conservative finite-sample check for
the zero-failure boundary case.

Power planning was performed on the development blocks before Campaign
C. The held-out run therefore used exactly ten independent blocks, four
episodes, and eight slots per episode. It was not stopped when a
threshold first appeared to pass, and no seed was replaced after
inspecting a result. The interrupted four-condition launch is kept
separate to make this stopping rule auditable.

The analysis distinguishes three kinds of zero. A zero outage can mean
that no failure was observed in the recorded slots; a zero work delta
can mean equal matched counts; and a missing physical timing value means
that no measurement exists. Only the first two are numeric simulator
outcomes. The third is represented as UNKNOWN in the audit and in the
latency figure, never as zero.

These estimands also constrain the prose. A condition-level PASS means
that the bound is below the registered threshold on that cell. It does
not imply a simultaneous population guarantee outside the factor grid. A
pooled mean is shown only for visualization and cannot override a
failing or unknown cell. This rule is applied to the periodic-baseline
cost failure and to the absent hardware trace.

\section{Protocol}\label{protocol}

\subsection{Frozen held-out Campaign
C}\label{frozen-held-out-campaign-c}

The corrected held-out campaign uses RIS sizes 16 and 32, user counts 2
and 4, error scales 0.12 and 0.28, and temporal correlations 0.70 and
0.98. Ten held-out seed blocks (1103, 1109, 1117, 1123, 1129, 1151,
1153, 1163, 1171, 1181) are disjoint from the four development blocks
(809, 811, 821, 823). Each condition-controller-seed stream contains
four episodes and eight slots. The four controllers are
posterior-triggered guard1 revision, periodic risk, always risk, and
periodic nominal reuse. This gives 640 streams and 20,480 rows. The
manifest classifies the run as \texttt{simulation\_only} and records the
source, configuration, and effective-configuration hashes.

The corrected run has zero missing or duplicate rows. Raw JSONL, JSON
summary, and CSV summary recomputation agree, and all paired comparisons
use matched seed, episode, and slot streams. These checks establish
reproducibility of the declared simulator; they do not turn simulator
rows into independent physical experiments.

\begin{figure*}[t]
\centering
\includegraphics[width=\textwidth]{../results/figures/campaign_c_20261007/figure01_protocol_provenance.pdf}
\caption{Protocol and provenance for the normalized RIS-only simulation. The panels enumerate the 16 error/rho and system-size conditions, the four matched controllers, the 20,480 scheduled rows, and the source/configuration hashes used for Campaign C. The \texttt{simulation\_only} label is a scope boundary: the diagram contains no measured RF waveform, phase-write completion, or physical slot timing.}\label{fig:campaign-01}
\end{figure*}

\subsection{Statistical unit and paired
analysis}\label{statistical-unit-and-paired-analysis}

The independent analysis unit is a complete seed block. A block contains
all episodes and slots generated from one seed. Paired deltas are formed
within a condition and seed block before resampling. The analysis uses
seed 20261031 and 10,000 bootstrap draws. Slots remain nested
observations and are never treated as independent channel replicates.
The same rule is used for outage, goodput, rate evaluations, refresh
frequency, and timing diagnostics.

\subsection{Absolute service gate}\label{absolute-service-gate}

Every scheduled slot remains in the denominator. A refusal, timeout,
missing decision, failed delivery, or failed phase write is one service
violation and zero delivered payload. The corrected Campaign C records
the resulting all-slot weak-user service metrics for every condition.
The maximum one-sided 95\% seed-block upper bound is 0.000, below the
frozen 0.05 threshold. This is an achieved result for the declared
normalized simulator, not a deployed service guarantee.

The corresponding all-slot service result is reported in Figure 4.

\subsection{Physical latency
boundary}\label{physical-latency-boundary}

The physical test would require a testbed-provided slot duration,
monotonic timestamps from the first available input through phase-write
completion, component decomposition, transfer timing, and an overrun
count. The present Jetson run has none of these physical records. Its
total-latency P95 is a simulator diagnostic and is reported without a
\texttt{0.10*T\_slot} threshold. The physical timing status is therefore
UNKNOWN and no hardware real-time claim is made.

The corresponding latency diagnostic is reported in Figure 5.

\subsection{Cost and baseline
freeze}\label{cost-and-baseline-freeze}

The ledger counts scenario construction, candidate search, validation,
and phase-control rate evaluations under the same accounting boundary.
The development selection rule chose the lowest-cost candidate
satisfying every primary development gate; the manifest was frozen
before the ten held-out blocks were read. The primary cost reference is
always-risk, as declared in the freeze file. Periodic-risk and
periodic-nominal comparisons are retained for context, but their
different refresh schedules make them unsuitable for the same 20\%
lower-bound cost assertion.

The cost and freeze outputs are reported in Figures 6 and 7.

\subsection{Independent held-out error/rho
grid}\label{independent-held-out-errorrho-grid}

Campaign C is the independent held-out error/rho grid. Its policy,
thresholds, seed list, candidate budget, and stopping rule were frozen
before evaluation. The resulting conditional comparison is valid for the
declared grid and matched simulator streams. It is not evidence of
universal distribution-shift robustness, mobility performance, or
calibrated uncertainty coverage outside the specified error scales and
temporal correlations.

The preserved campaign boundary is reported in Figure 8.

\section{Results}\label{results}

\subsection{Campaign integrity and
provenance}\label{campaign-integrity-and-provenance}

Campaign C completed on the Jetson with 20,480 of 20,480 planned rows.
The remote manifest and progress record both state \texttt{complete}; a
read-only process inventory found no active runner or verifier after
completion. The verifier reports 16 factor conditions, 64 summaries, 640
paired streams, and agreement between raw, JSON, and CSV recomputation.
The source SHA-256 is
\texttt{bb773999\allowbreak{}4c6cc109\allowbreak{}8ff43cdc\allowbreak{}a7aa39ff\allowbreak{}1c81ee00\allowbreak{}5707a185\allowbreak{}293b4358\allowbreak{}8957f4bd};
the configuration and effective-configuration hashes are recorded in the
manifest and completion receipt. The first wrong-subnet launch remains a
separate 2,450/5,120-row \texttt{running} artifact with no active
process and is explicitly excluded.

These records establish an auditable normalized simulation artifact.
They do not measure OTA propagation, payload errors on a physical radio,
phase-write completion, or a physical slot.

\subsection{Conditional controller
comparison}\label{conditional-controller-comparison}

Against always-risk, the posterior-triggered guard1 revision has a
maximum one-sided outage-delta upper bound of 0.000, a minimum goodput
relative lower 05 bound of −0.00615, and a minimum phase-control
rate-evaluation reduction lower 05 bound of 0.33398. The absolute
weak-user outage upper bound is 0.000 over the 16 conditions. Thus the
four frozen statistical gates pass on the declared Campaign C grid.

The result is a conditional trade-off, not dominance. Against
periodic-risk and periodic-nominal, service guardrails pass but the
work-reduction lower bound is negative in at least one condition. Those
baselines refresh on different schedules, so their work results are
retained as sensitivity evidence and are not used to revise the primary
always-risk gate after the held-out data.

\begin{figure*}[t]
\centering
\includegraphics[width=\textwidth]{../results/figures/campaign_c_20261007/figure02_paired_forest_always_risk.pdf}
\caption{Condition-level paired differences between posterior-triggered control and the always-risk reference. Each point is a condition-level estimate and each interval is obtained by resampling complete independent seed blocks (10,000 paired bootstrap draws); slots and episodes remain nested inside a block. The dashed lines are the preregistered outage, goodput, and rate-evaluation gates, not fitted confidence targets.}\label{fig:campaign-02}
\end{figure*}

\subsection{Error and temporal-factor
patterns}\label{error-and-temporal-factor-patterns}

The 16 condition cells separate error scale, temporal correlation, RIS
size, and user count. The heatmap reports the posterior controller's
phase-control rate evaluations, while the paired file reports the
matched reductions and goodput deltas. This separation shows where reuse
changes online work without collapsing the factor grid into a favorable
pooled mean. The pattern is a description of the normalized simulator;
it is not a causal estimate of a physical mobility process and does not
license interpolation beyond the declared error/rho cells.

\begin{figure}[t]
\centering
\includegraphics[width=\columnwidth]{../results/figures/campaign_c_20261007/figure03_error_rho_heatmap.pdf}
\caption{Posterior-triggered phase-control rate evaluations across the declared CSI-error scale and temporal-correlation grid. Each cell is the recorded online work count in the normalized simulator for one system-size and user-count condition; the annotations are counts, and the colorbar is a shared scale across all panels. Truth evaluation and physical phase writes are outside this online-work metric.}\label{fig:campaign-03}
\end{figure}

\subsection{Absolute weak-user
service}\label{absolute-weak-user-service}

The absolute service gate is closed for Campaign C. Every scheduled slot
is counted, and the maximum one-sided upper95 weak-user outage for the
posterior-triggered controller is 0.000. The service-violation field
remains available in the raw rows; its highest condition-level one-sided
95\% upper bound is 0.009375. As a conservative finite-sample
sensitivity, the attached Clopper--Pearson audit gives a maximum
weak-user upper bound of 0.009318 when the zero-failure condition cells
are treated as 320 scheduled weak-user slots; this audit does not
replace the registered seed-block bootstrap. The accounting includes any
refused or failed delivery rather than conditioning on admitted actions.
The safe statement is that the declared simulator campaign passed the
bound; it is not that a physical communication system is guaranteed to
meet 5\% outage.

\begin{figure}[t]
\centering
\includegraphics[width=\columnwidth]{../results/figures/campaign_c_20261007/figure04_absolute_outage_gate.pdf}
\caption{All-slot weak-user outage and service-violation diagnostics on the declared 16-condition Campaign C grid. Points are condition-level rates; whiskers are the registered one-sided seed-block upper95 bounds, with every scheduled slot retained in the denominator. The 0.05 line is the frozen absolute-outage gate, and a refusal, timeout, missing decision, or failed delivery is counted as a service failure rather than removed.}\label{fig:campaign-04}
\end{figure}

\subsection{Physical decision
latency}\label{physical-decision-latency}

The maximum posterior-triggered simulated total-latency P95 is 388.083
ms. This value includes the in-controller simulator timing scope
recorded in the manifest, but it does not include a physical RIS
transfer or phase-write completion. Because no physical \texttt{T\_slot}
or phase-commit trace exists, the latency gate is UNKNOWN. The paper
makes no statement about 10\%-of-slot feasibility, hardware real time,
or OTA operation.

\begin{figure}[t]
\centering
\includegraphics[width=\columnwidth]{../results/figures/campaign_c_20261007/figure05_latency_unknown_boundary.pdf}
\caption{Simulator timing components and the physical-evidence boundary. The plotted distributions are diagnostic wall-clock measurements from the Jetson execution path, summarized over the declared condition and controller records. They do not include a measured RIS phase-write completion or a physical slot duration; consequently no hardware-real-time threshold is drawn and the physical timing status remains UNKNOWN.}\label{fig:campaign-05}
\end{figure}

\subsection{Online work and conditional Pareto
interpretation}\label{online-work-and-conditional-pareto-interpretation}

The primary comparison shows a 0.33398 lower05 reduction in
phase-control rate evaluations versus always-risk while staying within
the declared outage and goodput bounds. This is a matched phase-control
evaluation result, not a total deployment-cost claim: physical transfer,
synchronization, and phase-write costs are unmeasured. The periodic
baseline failures illustrate why the cost denominator must be fixed
before interpreting reuse savings.

\begin{figure}[t]
\centering
\includegraphics[width=\columnwidth]{../results/figures/campaign_c_20261007/figure06_cost_service_conditional.pdf}
\caption{Conditional online-work and service view for the four matched controllers. Points and intervals use the same seed-block uncertainty unit as Figure 2. Always-risk is the preregistered primary cost reference; periodic risk and periodic nominal reuse are retained as sensitivity comparisons, whose different refresh schedules leave their all-condition work gate failed.}\label{fig:campaign-06}
\end{figure}

\subsection{Baseline fairness and tuning
freeze}\label{baseline-fairness-and-tuning-freeze}

The development analysis compared guard1, guard2, and adaptive
candidates on four development conditions. Guard1 revision was selected
because it passed the frozen service gates and had the largest work
reduction among passing candidates. The held-out seeds were disjoint and
the freeze file records the selection rule, thresholds, and hashes
before Campaign C. No held-out result was used to change the policy or
threshold.

\begin{figure}[t]
\centering
\includegraphics[width=\columnwidth]{../results/figures/campaign_c_20261007/figure07_freeze_protocol.pdf}
\caption{Freeze protocol separating development selection from held-out evaluation. Guard-1 revision selection uses disjoint development blocks; the held-out Campaign C blocks are read only after the policy, thresholds, seed lists, and source/configuration hashes are frozen. The ledger records all rows, failures, and analysis seeds needed to reproduce the reported intervals.}\label{fig:campaign-07}
\end{figure}

\subsection{Independent held-out error/rho
result}\label{independent-held-out-errorrho-result}

Campaign C supplies the independent held-out result that was missing
from the earlier draft. It supports the four conditional simulator
statements in Section 5.2. It does not support a claim of generalization
to arbitrary error distributions, mobility, hardware channels, or unseen
phase interfaces. The first partial launch is shown only as a preserved
failure record.

\begin{figure*}[t]
\centering
\includegraphics[width=\textwidth]{../results/figures/campaign_c_20261007/figure08_campaign_boundary.pdf}
\caption{Preserved Campaign A/B history and the corrected held-out Campaign C release boundary. A and B remain historical records, while C is the 20,480-row frozen evaluation used for the primary simulator gates. The periodic-baseline work failure and the absence of physical phase-write/slot timing remain visible as evidence boundaries rather than being relabelled as successes.}\label{fig:campaign-08}
\end{figure*}

\subsection{Tables and reproducibility
records}\label{tables-and-reproducibility-records}

Table 1 defines the normalized link and CSI assumptions. Table 2 defines
all-slot service accounting. Table 3 separates simulator timing from
physical timing. Table 4 lists the frozen configuration and source
hashes. Table 5 reports seed-block paired intervals. Table 6 is the
release ledger, with physical phase-write timing marked UNKNOWN rather
than imputed from simulator wall time. All figures are generated from
the CSV/JSON artifacts and have the provenance manifest
\path{results/figures/campaign_c_20261007/figure_provenance.json}.
The exact zero-failure audit is stored at
\path{results/controller_shift_test_final_guard1_revision_full_20261007_jetson_analysis/exact_outage_sensitivity.json}.


\begin{table}[t]
\caption{Normalized simulator and CSI settings.}
\label{tab:settings}
\scriptsize
\begin{tabularx}{\columnwidth}{@{}l l X@{}}
\toprule
Quantity & Frozen setting & Evidence boundary \\
\midrule
RIS elements & 16, 32 & normalized simulator input \\
Users & 2, 4 & matched multi-user streams \\
Phase alphabet & 2-bit & discrete action space \\
Error scale & 0.12, 0.28 & simulated estimation error \\
Temporal correlation & 0.70, 0.98 & AR(1) simulator factor \\
Scenario ensemble & 4; guard 1 & decision surrogate, not calibrated posterior \\
Direct BS--user path & omitted & no physical channel claim \\
\bottomrule
\end{tabularx}
\end{table}

\begin{table}[t]
\caption{All-slot service and online-work accounting.}
\label{tab:accounting}
\scriptsize
\begin{tabularx}{\columnwidth}{@{}l X@{}}
\toprule
Record & Rule \\
\midrule
Scheduled slot & always remains in the denominator \\
Refusal, timeout or failed delivery & zero payload and one service violation \\
Weak-user outage & condition-level reference user, all scheduled slots \\
Primary work metric & phase-control rate evaluations \\
Truth evaluation & scores an action; excluded from online controller work \\
Initialization & recorded in campaign totals, not silently amortized \\
\bottomrule
\end{tabularx}
\end{table}

\begin{table}[t]
\caption{Timing evidence boundary.}
\label{tab:timing}
\scriptsize
\begin{tabularx}{\columnwidth}{@{}l X@{}}
\toprule
Quantity & Status in Campaign C \\
\midrule
Estimator/search/validation timing & simulator diagnostic \\
Transfer and synchronization & not a physical measurement \\
RIS phase-write completion & UNKNOWN; no endpoint detected \\
Physical $T_{\mathrm{slot}}$ and overruns & UNKNOWN; no slot trace \\
\bottomrule
\end{tabularx}
\end{table}

\begin{table}[t]
\caption{Frozen Campaign C protocol and hashes.}
\label{tab:protocol}
\scriptsize
\begin{tabularx}{\columnwidth}{@{}l X@{}}
\toprule
Item & Value \\
\midrule
Rows / conditions / streams & 20,480 / 16 / 640 \\
Seed blocks & 10 held-out; 4 development \\
Episodes / slots & 4 / 8 per stream \\
Bootstrap & 10,000 draws; seed 20261031 \\
Source SHA-256 & \texttt{bb773999...57f4bd} \\
Configuration SHA-256 & \texttt{903fc0a2...f35e2b} \\
Effective-config SHA-256 & \texttt{60f9fc84...c2b4530} \\
\bottomrule
\end{tabularx}
\end{table}

\begin{table}[t]
\caption{Primary paired Campaign C gate results versus always-risk.}
\label{tab:gates}
\scriptsize
\begin{tabularx}{\columnwidth}{@{}p{0.28\columnwidth} X p{0.14\columnwidth}@{}}
\toprule
Estimand & Bound & Status \\
\midrule
Weak-user outage & upper95 $=0.000$; target $\leq0.05$ & PASS \\
Outage delta & upper95 $=0.000$; target $\leq0.01$ & PASS \\
Goodput relative delta & lower05 $=-0.00615$; target $\geq-0.05$ & PASS \\
Rate-evaluation reduction & lower05 $=0.33398$; target $\geq0.20$ & PASS \\
\bottomrule
\end{tabularx}
\end{table}

\begin{table}[t]
\caption{Release ledger and evidence boundaries.}
\label{tab:release}
\scriptsize
\begin{tabularx}{\columnwidth}{@{}l X@{}}
\toprule
Item & Release status \\
\midrule
Always-risk statistical gates & PASS on declared normalized grid \\
Periodic-risk / nominal work gate & FAIL; sensitivity only \\
Exact zero-failure outage audit & diagnostic; max upper95 $=0.009318$ \\
Physical phase-write and slot timing & UNKNOWN / NO-GO \\
OTA and deployment feasibility & UNKNOWN / NO-GO \\
Learned-policy superiority & not evaluated or claimed \\
\bottomrule
\end{tabularx}
\end{table}

\section{Discussion}\label{discussion}

\subsection{Interpretation within the simulation
scope}\label{interpretation-within-the-simulation-scope}

The completed evidence supports an inspectable causal controller and a
condition-specific reliability--computation trade-off. The scenario
ensemble makes the uncertainty assumption visible, the lower-tail check
makes the weak-user criterion explicit, and the trigger determines when
refresh work is spent. The held-out analysis shows that this design can
preserve the always-risk service reference on the declared normalized
grid while reducing phase-control rate evaluations. It does not show
that the controller is globally optimal or universally better.

\subsection{Reliability accounting}\label{reliability-accounting}

The all-slot denominator prevents a controller from appearing reliable
by refusing difficult actions. Campaign C passes the absolute weak-user
upper95 bound and the relative outage/goodput guardrails under this
definition. The interpretation remains conditional because the channel,
waveform, payload, and error model are simulated. A deployed service
guarantee would require physical channel and payload observations under
a separately registered sampling design.

\subsection{Decision budget and physical
realizability}\label{decision-budget-and-physical-realizability}

The simulated latency diagnostic is useful for comparing code paths, but
it is not a physical slot or a phase-write measurement. The absence of a
RIS control board or serial endpoint is recorded as an evidence
limitation rather than filled with a nominal interface assumption. A
future physical test must expose the exact slot duration, monotonic
timestamps, transfer and commit intervals, overruns, and payload/error
outcomes before any real-time claim is made.

\subsection{Independent shifts and baseline
fairness}\label{independent-shifts-and-baseline-fairness}

The held-out error/rho grid closes the registered conditional
comparison. The result is still bounded by its factors and seed blocks.
The periodic baseline work failures demonstrate that refresh schedules
affect the denominator; this is why the always-risk comparison was
frozen as the primary cost reference and the periodic results remain
visible rather than tuned away.

\subsection{Relation to prior work and
novelty}\label{relation-to-prior-work-and-novelty}

Robust RIS beamforming and outage-constrained design are established
directions {[}7,8,21,28,29,33,34{]}. Foundational RIS models and surveys
frame the propagation and control abstractions {[}2,6,11,13{]}, while
temporal estimation, aging, and prediction have been studied with
Kalman, two-timescale, and predictive mechanisms
{[}14,22--24,26,27,30--32{]}. Event-triggered estimation, deep
reinforcement learning, meta-learning, and closed-loop control also have
nearby precedents {[}5,9,10,12,19,25{]}. The candidate gap is the
narrower systems intersection of an inspectable posterior-like error
ensemble, finite lower-tail feasibility, event-triggered phase reuse,
and an audited reliability/rate/refresh/work ledger under a frozen
held-out protocol. Additional practical RIS, traffic-reliability, ISAC,
codebook, tracking, and secure STAR-RIS studies in the ledger delimit
the neighboring application spaces {[}1,3,4,15--18,20{]}. These
citations establish neighboring problem families; metadata-only entries
do not support detailed algorithm claims. The present study does not
train or evaluate a deep or reinforcement-learning policy, so it makes
no claim of superiority over learned control. The paper must not claim
``first,'' ``novel,'' DRL, mobility, or physical deployment based on
this campaign.

\subsection{Limitations and next
tests}\label{limitations-and-next-tests}

The normalized simulator omits direct BS--user propagation, absolute
path loss, waveform and payload measurement, calibrated hardware timing,
and a physical RIS controller. The finite scenario ensemble is a design
surrogate and can miss error modes outside its inflation model. Two-bit
actions and finite candidate budgets constrain the search. The weak-user
metric is a condition-level reference metric defined from the matched
always-risk trace; it should not be interpreted as a fixed user identity
or as a population-level worst-user guarantee. The cost result covers
phase-control rate evaluations; transfer, synchronization, and
phase-write costs remain unmeasured. The next clean extension is a
physical interface qualification with a registered slot trace, followed
by a new held-out hardware campaign; it must not overwrite Campaign C or
the preserved partial record.

\subsection{Availability and audit
trail}\label{availability-and-audit-trail}

The source implementation, YAML configurations, unit tests, raw JSONL
rows, summary CSV files, paired bootstrap records, gate audit, and
figure provenance are retained in the project workspace. The corrected
Campaign C completion record links every reported number to source,
configuration, effective configuration, manifest, summary, and raw-file
SHA-256 values. The preserved interrupted launch is stored beside the
completed run rather than deleted, so an independent reader can
distinguish a failed start from the final evidence. The local verifier
reproduces row uniqueness, condition completeness, and raw/JSON/CSV
agreement. Runtime replay is treated separately: the campaign
environment is Jetson Python 3.12.3 with NumPy 2.5.3, whereas a Windows
replay uses a different runtime and is not claimed as a bitwise
execution match. This separation keeps a structural audit from being
mistaken for a hardware or software-environment replication result.

\section{Conclusion}\label{conclusion}

We evaluated a transparent causal RIS phase-reuse controller for
temporally drifting CSI estimation error. The controller combines a
finite error-scenario ensemble, a per-user lower-tail feasibility check,
a quantized candidate search, and an explicit refresh/reuse trigger. In
the corrected 20,480-row Campaign C, all four frozen statistical gates
pass against always-risk on the declared 16-condition normalized
simulator grid: the absolute weak-user outage upper95 is 0.000, the
outage-delta upper95 is 0.000, the goodput lower05 is −0.00615, and the
rate-evaluation reduction lower05 is 0.33398.

The result is a conditional simulation finding. Periodic baseline work
gates remain failed and visible, the first wrong-subnet launch remains
excluded, and physical phase-write completion, slot timing, OTA
performance, and universal generalization remain unverified. A
submission can therefore be framed as a reproducible normalized RIS
control and evaluation study, but it must not claim hardware real time,
deployment, calibrated chance constraints, learned control, or universal
superiority.

\section*{Submission statements}\label{submission-statements}

\subsection*{Data and code availability}\label{data-and-code-availability}

The source implementation, versioned YAML configurations, unit tests,
raw Campaign C JSONL rows, summaries, paired bootstrap indices, gate
audits, and figure-generation scripts are retained in the project
workspace. The hash-linked completion record is
\path{results/controller_shift_test_final_guard1_revision_full_20261007_jetson_analysis/campaign_c_completion.json}.
The raw rows are simulator-generated rather than measurements from human
participants or a physical radio. The reproducibility archive is
publicly available at
https://github.com/pengliangtian06-commits/transparent-risk-sensitive-ris-control
and archived at Zenodo (DOI: 10.5281/zenodo.23220663).

\subsection*{Ethics and competing interests}\label{ethics-and-competing-interests}

This study uses no human participants, animals, personal data, or
clinical records; an ethics review is therefore not applicable to the
reported simulation. The authors declare no competing interests.

\subsection*{Funding}\label{funding}

The authors received no external funding for this work.

\subsection*{Author contributions}\label{author-contributions}

Pengliang Tian led the technical work and writing of the original draft.
Pengyuan Zhang handled manuscript typesetting and formatting. Haoshan
Sun checked and verified the data. Changjiang Zhang supervised the work
and reviewed and edited the manuscript. All authors reviewed and
approved the submitted version.

\subsection*{AI-assisted work disclosure}\label{ai-assisted-work-disclosure}

An AI assistant was used for code inspection, reproducibility auditing,
figure generation scripts, and language drafting. The authors must
verify all technical claims, citations, and final wording and should
align this statement with the target journal's current Elsevier/Physical
Communication policy.

\section*{References}\label{references-and-verification-status}

The following 34 records form the current manuscript candidate pool.
Bibliographic identity was checked through Crossref; \texttt{FT} marks a
local full-text reading, and \texttt{MD} marks metadata/abstract-level
screening. Detailed claims use only records whose evidence level
supports them. The final submission should select 30--40 entries after
the last full-text screen, DOI deduplication, publisher-version check,
and current CAS/SCI classification check.

{[}1{]} Jialing Xu, Jianbin Xue, Han Zhang, Xiangrui Guan. Robust secure
transmission for STAR-RIS assisted ISAC against collaborative
eavesdropper with imperfect CSI. Physical Communication, 2026. DOI:
10.1016/j.phycom.2026.103011. (MD) {[}2{]} Ertugrul Basar, Marco Di
Renzo, Julien De Rosny, Merouane Debbah, Mohamed-Slim Alouini, Rui
Zhang. Wireless Communications Through Reconfigurable Intelligent
Surfaces. IEEE Access, 2019. DOI: 10.1109/access.2019.2935192. (MD)
{[}3{]} Hamidreza Khaleghi, Stéphane Paquelet. Adaptive Low-Overhead
Channel Estimation Tracking in RIS-Assisted Systems. IEEE Access, 2025.
DOI: 10.1109/ACCESS.2025.3570772. (MD) {[}4{]} Han Han, Tao Jiang, Wei
Yu. Active Beam Tracking with Reconfigurable Intelligent Surface. ICASSP
2023 - 2023 IEEE International Conference on Acoustics, Speech and
Signal Processing (ICASSP), 2023. DOI:
10.1109/ICASSP49357.2023.10097134. (MD) {[}5{]} Ramin Hashemi, Samad
Ali, Nurul Huda Mahmood, Matti Latva-Aho. Deep Reinforcement Learning
for Practical Phase-Shift Optimization in RIS-Aided MISO URLLC Systems.
IEEE Internet of Things Journal, 2023. DOI: 10.1109/JIOT.2022.3232962.
(MD) {[}6{]} Marco Di Renzo, Alessio Zappone, Merouane Debbah,
Mohamed-Slim Alouini, Chau Yuen, Julien de Rosny, et al.. Smart Radio
Environments Empowered by Reconfigurable Intelligent Surfaces: How It
Works, State of Research, and The Road Ahead. IEEE Journal on Selected
Areas in Communications, 2020. DOI: 10.1109/jsac.2020.3007211. (MD)
{[}7{]} Gui Zhou, Cunhua Pan, Hong Ren, Kezhi Wang, Marco Di Renzo,
Arumugam Nallanathan. Robust Beamforming Design for Intelligent
Reflecting Surface Aided MISO Communication Systems. IEEE Wireless
Communications Letters, 2020. DOI: 10.1109/LWC.2020.3000490. (MD)
{[}8{]} Hui Gao, Kai Cui, Chongwen Huang, Chau Yuen. Robust Beamforming
for RIS-Assisted Wireless Communications With Discrete Phase Shifts.
IEEE Wireless Communications Letters, 2021. DOI:
10.1109/lwc.2021.3107319. (MD) {[}9{]} Xinquan Wang, Fenghao Zhu,
Chongwen Huang, Ahmed Alhammadi, Faouzi Bader, Zhaoyang Zhang, et al..
Robust Beamforming With Gradient-Based Liquid Neural Network. IEEE
Wireless Communications Letters, 2024. DOI: 10.1109/LWC.2024.3436576.
(MD) {[}10{]} Luis Carlos Mathias, Rafael Marasca Martins, Taufik Abrão.
A Closed-Loop Model Predictive Control for Beamforming in RIS-Aided
Communications. IEEE Wireless Communications Letters, 2026. DOI:
10.1109/LWC.2026.3701340. (MD) {[}11{]} Xiaojun Yuan, Ying-Jun Angela
Zhang, Yuanming Shi, Wenjing Yan, Hang Liu.
Reconfigurable-Intelligent-Surface Empowered Wireless Communications:
Challenges and Opportunities. IEEE Wireless Communications, 2021. DOI:
10.1109/mwc.001.2000256. (MD) {[}12{]} Dwina Agustin Putri, Irma Zakia,
Hendrawan Hendrawan, Robithoh Annur. Joint Curriculum-Enhanced DRL and
Event-Triggered Channel Re-Estimation for Hybrid Beamforming in Dynamic
mmWave Environments. IEEE Open Journal of the Communications Society,
2025. DOI: 10.1109/OJCOMS.2025.3642009. (MD) {[}13{]} Qingqing Wu,
Shuowen Zhang, Beixiong Zheng, Changsheng You, Rui Zhang. Intelligent
Reflecting Surface-Aided Wireless Communications: A Tutorial. IEEE
Transactions on Communications, 2021. DOI: 10.1109/tcomm.2021.3051897.
(MD) {[}14{]} Chen Hu, Linglong Dai, Shuangfeng Han, Xiaoyun Wang.
Two-Timescale Channel Estimation for Reconfigurable Intelligent Surface
Aided Wireless Communications. IEEE Transactions on Communications,
2021. DOI: 10.1109/tcomm.2021.3072729. (MD) {[}15{]} Mohammed
Almekhlafi, Mohamed Amine Arfaoui, Mohamed Elhattab, Chadi Assi, Ali
Ghrayeb. Joint Resource Allocation and Phase Shift Optimization for
RIS-Aided eMBB/URLLC Traffic Multiplexing. IEEE Transactions on
Communications, 2022. DOI: 10.1109/TCOMM.2021.3127265. (MD) {[}16{]}
Kangda Zhi, Cunhua Pan, Hong Ren, Kezhi Wang. Power Scaling Law Analysis
and Phase Shift Optimization of RIS-Aided Massive MIMO Systems With
Statistical CSI. IEEE Transactions on Communications, 2022. DOI:
10.1109/TCOMM.2022.3162580. (MD) {[}17{]} Mohamed Rihan, Alessio
Zappone, Stefano Buzzi. Robust RIS-Assisted MIMO Communication-Radar
Coexistence: Joint Beamforming and Waveform Design. IEEE Transactions on
Communications, 2023. DOI: 10.1109/TCOMM.2023.3298983. (MD) {[}18{]}
Yongqing Xu, Yong Li, J. Andrew Zhang, Marco Di Renzo, Tony Q. S. Quek.
Joint Beamforming for RIS-Assisted Integrated Sensing and Communication
Systems. IEEE Transactions on Communications, 2024. DOI:
10.1109/TCOMM.2023.3344143. (MD) {[}19{]} Min Wu, Kefeng Guo, Xingwang
Li, Zhi Lin, Yongpeng Wu, Theodoros A. Tsiftsis, et al.. Deep
Reinforcement Learning-Based Energy Efficiency Optimization for
RIS-Aided Integrated Satellite-Aerial-Terrestrial Relay Networks. IEEE
Transactions on Communications, 2024. DOI: 10.1109/TCOMM.2024.3370618.
(MD) {[}20{]} Zhiheng Yu, Jiancheng An, Ertugrul Basar, Lu Gan, Chau
Yuen. Environment-Aware Codebook Design for RIS-Assisted MU-MISO
Communications: Implementation and Performance Analysis. IEEE
Transactions on Communications, 2024. DOI: 10.1109/TCOMM.2024.3415619.
(MD) {[}21{]} Ming-Min Zhao, An Liu, Rui Zhang. Outage-Constrained
Robust Beamforming for Intelligent Reflecting Surface Aided Wireless
Communication. IEEE Transactions on Signal Processing, 2021. DOI:
10.1109/TSP.2021.3056899. (FT) {[}22{]} Anastasios Papazafeiropoulos,
Ioannis Krikidis, Pandelis Kourtessis. Impact of Channel Aging on
Reconfigurable Intelligent Surface Aided Massive MIMO Systems With
Statistical CSI. IEEE Transactions on Vehicular Technology, 2023. DOI:
10.1109/TVT.2022.3203796. (MD) {[}23{]} Yan Zhang, Jiayi Zhang, Huahua
Xiao, Derrick Wing Kwan Ng, Bo Ai. Channel Aging-Aware Precoding for
RIS-Aided Multi-User Communications. IEEE Transactions on Vehicular
Technology, 2023. DOI: 10.1109/TVT.2022.3210548. (MD) {[}24{]} Enyu Shi,
Jiayi Zhang, Jiakang Zheng, Bo Ai, Derrick Wing Kwan Ng. RIS-Aided
Cell-Free Massive MIMO Systems With Channel Aging. IEEE Transactions on
Vehicular Technology, 2024. DOI: 10.1109/TVT.2024.3376557. (MD) {[}25{]}
Yiyang Zhu, Enyu Shi, Ziheng Liu, Jiayi Zhang, Bo Ai. Multi-Agent
Reinforcement Learning-Based Joint Precoding and Phase Shift
Optimization for RIS-Aided Cell-Free Massive MIMO Systems. IEEE
Transactions on Vehicular Technology, 2024. DOI:
10.1109/TVT.2024.3392883. (MD) {[}26{]} Nipuni Ginige, Arthur Sousa de
Sena, Nurul Huda Mahmood, Nandana Rajatheva, Matti Latva-Aho. Efficient
Channel Prediction for Beyond Diagonal RIS-Assisted MIMO Systems With
Channel Aging. IEEE Transactions on Vehicular Technology, 2025. DOI:
10.1109/TVT.2025.3556277. (MD) {[}27{]} Yi Wei, Ming-Min Zhao, An Liu,
Min-Jian Zhao. Channel Tracking and Prediction for IRS-Aided Wireless
Communications. IEEE Transactions on Wireless Communications, 2023. DOI:
10.1109/TWC.2022.3196313. (MD) {[}28{]} Zhen Chen, Jie Tang, Xiu Yin
Zhang, Qingqing Wu, Gaojie Chen, Kai-Kit Wong. Robust Hybrid Beamforming
Design for Multi-RIS Assisted MIMO System With Imperfect CSI. IEEE
Transactions on Wireless Communications, 2023. DOI:
10.1109/TWC.2022.3222218. (FT) {[}29{]} Yasaman Omid, S. M. Mahdi
Shahabi, Cunhua Pan, Yansha Deng, Arumugam Nallanathan. Robust
Beamforming Design for an IRS-Aided NOMA Communication System With CSI
Uncertainty. IEEE Transactions on Wireless Communications, 2024. DOI:
10.1109/TWC.2023.3283140. (MD) {[}30{]} Danyang Yu, Gan Zheng, Arman
Shojaeifard, Sangarapillai Lambotharan, Yi Liu. Kalman Filter Based
Channel Tracking for RIS-Assisted Multi-User Networks. IEEE Transactions
on Wireless Communications, 2024. DOI: 10.1109/TWC.2023.3312426. (FT)
{[}31{]} Fanghao Xia, Zesong Fei, Jingxuan Huang, Xinyi Wang, Ruixiang
Wang, Weijie Yuan, et al.. Sensing-Enabled Predictive Beamforming Design
for RIS-Assisted V2I Systems: A Deep Learning Approach. IEEE
Transactions on Wireless Communications, 2024. DOI:
10.1109/TWC.2023.3327362. (MD) {[}32{]} Qian Zhang, Ju Liu, Haoge Tang,
Zheng Dong, Yonghui Li. Practical RIS-Aided Multiuser Communications
With Imperfect CSI: Practical Model, Amplitude Feedback, and Beamforming
Optimization. IEEE Transactions on Wireless Communications, 2024. DOI:
10.1109/twc.2024.3427695. (MD) {[}33{]} Fenghao Zhu, Xinquan Wang,
Chongwen Huang, Zhaohui Yang, Xiaoming Chen, Ahmed Al Hammadi, et al..
Robust Beamforming for RIS-Aided Communications: Gradient-Based Manifold
Meta Learning. IEEE Transactions on Wireless Communications, 2024. DOI:
10.1109/TWC.2024.3435023. (MD) {[}34{]} Yan Yang, Shuping Dang, Miaowen
Wen, Bo Ai, Rose Qingyang Hu. Blockage-Aware Robust Beamforming in
RIS-Aided Mobile Millimeter Wave MIMO Systems. IEEE Transactions on
Wireless Communications, 2024. DOI: 10.1109/TWC.2024.3447895. (MD)

Metadata audit artifact:
\path{results/reference_metadata_final_20261007.json}; selection
ledger: \path{literature/FINAL_BIBLIOGRAPHY_METADATA_20261007.md}.
\end{document}
